%% Elements of Nested Fibrational Cosmology
%% Book IV: The Universal Source and Completion Ladder
%% Final-Pass Canonical Draft
%%
%% Status legend:
%%   [U]  Unconditional  — proved from standing definitions and prior Unc results
%%   [C]  Conditional    — depends on explicitly declared hypotheses
%%   [D]  Definition-structural — structural, no proof obligation
%%   [R]  Remark only    — expository, non-load-bearing
%%   [O]  Open obligation — not yet proved

\documentclass[12pt,oneside]{amsbook}

\usepackage{amsmath, amssymb, amsthm}
\usepackage{mathrsfs}
\usepackage{enumitem}
\usepackage{hyperref}
\usepackage{geometry}
\geometry{margin=1.2in}

%% ---------------------------------------------------------------
%% Theorem environments
%% ---------------------------------------------------------------

\theoremstyle{plain}
\newtheorem{theorem}{Theorem}[section]
\newtheorem{lemma}[theorem]{Lemma}
\newtheorem{proposition}[theorem]{Proposition}
\newtheorem{corollary}[theorem]{Corollary}

\theoremstyle{definition}
\newtheorem{definition}[theorem]{Definition}
\newtheorem{standingrule}[theorem]{Standing Rule}

\theoremstyle{remark}
\newtheorem{remark}[theorem]{Remark}
\newtheorem{scholium}[theorem]{Scholium}

%% Status tag macro
\newcommand{\status}[1]{\textup{[#1]}}

%% ---------------------------------------------------------------
%% Notation
%% ---------------------------------------------------------------
%%  From Books I--III (all freely used):
%%    \Omega, L, \Sigma_n, \Sigma_\infty, \delta, \Delta
%%    U, \partial_{\rho^*}U, O(U), C_{\rho^*}(U)
%%    C_n, \Theta_n, E_n, B_n  (Book III stability functional)
%%
%%  New in Book IV:
%%    \mathrm{Reg}         base category of admissible regimes
%%    \mathrm{POT}         category of persistent observational-transfer objects
%%    F_R                  fiber over regime R
%%    \mathbf{D}_{\mathrm{pers}}: \mathrm{Reg} \to \mathrm{POT}
%%                         persistence diagram functor
%%    \tau_\iota           coherent transition map over enlargement \iota
%%    I(R,U), I_{\mathrm{nl}}(R,U)  transport invariant (linear/nonlinear)
%%    K^*                  minimal carrier sub-object of U
%%    T^*                  minimal observable-sufficient target
%%    \pi: K^* \to T^*     canonical projection
%%    U                    universal source object (roman, not italic)

\begin{document}

\title{Book IV\\[6pt]The Universal Source and Completion Ladder}
\author{Elements of Nested Fibrational Cosmology}
\date{}
\maketitle

\tableofcontents

\bigskip

%% ---------------------------------------------------------------
\section*{IV.0\quad Purpose}
%% ---------------------------------------------------------------

Book~IV is the apex of the canonical spine.  Its task is to establish
the existence, certified role, and correct uniqueness notion of the
universal source object.

This book does not certify any specific physical endpoint.  It
certifies the upstream source from which all lawful descendants must
factor.  Its burden is fourfold:
\begin{enumerate}
  \item define the source arena (the POT category),
  \item state the completion ladder leading to the universal source,
  \item prove the existence and source-role of the universal source
        object $\mathbf{U}$, and
  \item mark the exact downstream boundary that remains after
        $\mathbf{U}$ is constructed.
\end{enumerate}

Books~I--III have established the machinery required for this
construction.  Book~I supplied the observational quotient and defect
bookkeeping.  Book~II supplied the boundary law and collar structure.
Book~III supplied the persistence and transport machinery, including
the Unified Coercive-Transport Inequality, and defined
observational-transfer objects as the comparison arena.  Book~IV
assembles all of this into a single universal object.

\begin{standingrule}[\status{D}]\label{sr:source-not-branch}
Book~IV is not yet about gauge sectors, spacetime structure, matter
representations, or branch endpoints.  It certifies the upstream
object from which such descendants, if lawful, must later factor.
Representation content is not primitive source content.  It enters
only through explicit realization structure in derived branch books.
\end{standingrule}

%% ---------------------------------------------------------------
\section{The Source Arena}
%% ---------------------------------------------------------------

\subsection*{The base category of admissible regimes}

\begin{definition}[\status{D}]\label{def:Reg-regime}
An \emph{admissible regime} is a pair $R = (\Omega_R, \mathcal{A}_R)$
consisting of a finite configuration space $\Omega_R$ together with a
collection $\mathcal{A}_R$ of admissible test families on $\Omega_R$
in the sense of Book~I.
\end{definition}

\begin{definition}[\status{D}]\label{def:Reg-enlargement}
An \emph{admissible enlargement} $\iota: R \hookrightarrow R'$ is an
injective map $\Omega_R \hookrightarrow \Omega_{R'}$ compatible with
admissible tests: the enlarged regime introduces no new primitive data
not already licensed by the upstream machinery.
\end{definition}

\begin{definition}[\status{D}]\label{def:Reg}
The \emph{base category} $\mathrm{Reg}$ has admissible regimes as
objects and admissible enlargements as morphisms.  Composition is
composition of injective enlargement maps.
\end{definition}

\begin{remark}[\status{R}]\label{rmk:Reg-not-poset}
$\mathrm{Reg}$ is not merely a poset.  Different enlargement paths
$R \hookrightarrow R' \hookrightarrow R''$ and $R \hookrightarrow R''
\hookrightarrow R'''$ may produce different intermediate regimes.  The
category structure records this path information.  This path-dependence
is exactly what the Overlap Route-Agreement Theorem (Section~\ref{sec:ladder})
must control before a universal object can be assembled.
\end{remark}

\subsection*{POT fibers}

\begin{definition}[\status{D}]\label{def:fiber}
The \emph{fiber} $F_R$ over an admissible regime $R$ consists of the
following licensed data, all inherited from Books~I--III:
\begin{enumerate}[label=\textup{(\arabic*)}]
  \item \emph{Stabilized observational quotient} $\Sigma_R$: the
        behavioral congruence quotient of $\Omega_R$ by the admissible
        tests $\mathcal{A}_R$ (Book~I, Theorem~I.6.1).
  \item \emph{Continuation-enlargement structure} $\Lambda_R$: the
        continuation data on $R$, including continuation histories,
        cumulative defect counts $\Lambda(h)$, and the defect-ledger
        structure (Book~I, Section~I.5).
  \item \emph{Transfer structure} $T_R$: the admissible transfer maps
        on $R$, including the boundary-law map $\Phi_R: C_{\rho^*}(R)
        \to O(R)$ under LS-2 (Book~II, Corollary~II.4.2) and the
        interface-relative stability structure (Book~III,
        Definitions~III.5.1--III.5.6).
  \item \emph{Representation layer} $\mathcal{R}_R$: initially empty
        for all $R$.  Populated only when a bridge theorem licenses it.
  \item \emph{Defect-ledger data} $D_R$: the defect ledger recording,
        for each continuation history $h$ on $R$, the cumulative
        defect count $\Lambda(h)$ and its decomposition (Book~I,
        Definition~I.5.3).
\end{enumerate}
\end{definition}

\begin{remark}[\status{R}]\label{rmk:rep-layer-empty}
Item~(4) in Definition~\ref{def:fiber} is the anti-smuggling firewall
at the categorical level.  The fiber is defined without assuming
sector persistence, canonical projection, or continuum realization.
Those features enter only when the corresponding bridge theorems are
proved.  At that point the representation layer of the relevant fiber
is populated by a downstream functor, not retroactively assumed.  No
content may be silently pre-loaded into the representation layer.
\end{remark}

\subsection*{Objects and morphisms of POT}

\begin{definition}[\status{D}]\label{def:POT-object}
A \emph{persistent observational-transfer object} (POT object) is a
pair $(R, F_R)$ where $R$ is an admissible regime and $F_R$ is its
fiber.
\end{definition}

\begin{definition}[\status{D}]\label{def:POT-morphism}
A \emph{POT-morphism} $f: X_R \to X'_{R'}$ over an admissible
enlargement $\iota: R \hookrightarrow R'$ consists of
structure-preserving maps between fiber components:
\begin{enumerate}[label=\textup{(\arabic*)}]
  \item $f_\Sigma: \Sigma_R \to \Sigma_{R'}$ compatible with
        observational quotient structure;
  \item $f_\Lambda: \Lambda_R \to \Lambda_{R'}$ compatible with
        continuation histories and defect-ledger data, satisfying
        $\Lambda_{R'}(f_\Lambda(h)) = \Lambda_R(h) + \delta_\iota(h)$
        for an explicit defect increment $\delta_\iota(h) \geq 0$;
  \item $f_T: T_R \to T_{R'}$ compatible with the boundary-law maps
        and the LS-2 collar-outcome discipline; and
  \item $f_D: D_R \to D_{R'}$ consistent with defect-ledger
        additivity.
\end{enumerate}
The representation layer, when non-empty, must also be transported
compatibly.
\end{definition}

\begin{proposition}[\status{U}]\label{prop:POT-compose}
POT-morphisms compose: if $f: X_R \to X'_{R'}$ lies over $\iota:
R \hookrightarrow R'$ and $g: X'_{R'} \to X''_{R''}$ lies over
$\kappa: R' \hookrightarrow R''$, then $g \circ f$ lies over $\kappa
\circ \iota$ and is a valid POT-morphism.
\end{proposition}

\begin{proof}
Composition of the four component maps is straightforward.
Defect-ledger additivity requires
\[
  \delta_{\kappa \circ \iota}(h)
  = \delta_\kappa(f_\Lambda(h)) + \delta_\iota(h),
\]
which holds by the additive structure of the continuation-enlargement
data (Book~I, Proposition~I.5.4).
\end{proof}

\begin{definition}[\status{D}]\label{def:POT-category}
The \emph{category of persistent observational-transfer objects}
$\mathrm{POT}$ has POT objects as objects, POT-morphisms as morphisms,
and forgetful functor $\pi: \mathrm{POT} \to \mathrm{Reg}$ sending
$(R, F_R) \mapsto R$.  $\mathrm{POT}$ is a Grothendieck fibration over
$\mathrm{Reg}$.
\end{definition}

\begin{remark}[\status{R}]\label{rmk:POT-not-presheaf}
POT is richer than a presheaf category on $\mathrm{Reg}$: the fibers
carry quotient structure, defect-ledger data, and boundary-law maps
from Books~I--III.  This extra structure is what makes the transport
invariant a POT-native quantity rather than an external assignment.
\end{remark}

\subsection*{The transport invariant}

\begin{definition}[\status{D}]\label{def:transport-invariant}
For $(R, F_R) \in \mathrm{POT}$ and a local region $U \subseteq
\Omega_R$, the \emph{compression-ledger transport invariant} is
\[
  I(R, U) \;=\; \log|O_R(U)| \;+\; \lambda\,\Lambda_R(U)
            \;+\; \mu\,\log|C_R(U)|,
\]
where $O_R(U)$ is the set of realized stabilized local outcomes,
$C_R(U) = C_{\rho^*}(U)$ is the set of realized stabilized collar
types under LS-2, $\Lambda_R(U)$ is the cumulative defect-ledger count
transported to $U$, and $\lambda, \mu \geq 0$ are normalization
coefficients fixed in Section~\ref{sec:ladder}.

The \emph{boundary-corrected nonlinear extension} is
\[
  I_{\mathrm{nl}}(R, U)
  \;=\; I(R,U) \;+\; \nu\,\Lambda_R(U) \cdot \log|C_R(U)|,
\]
with $\nu \geq 0$ an additional interaction coefficient.
\end{definition}

\begin{remark}[\status{R}]\label{rmk:invariant-POT-native}
The quantities $|O_R(U)|$ and $|C_R(U)|$ are defined at the level of
realized stabilized data, not raw token counts: they are certified by
the observational quotient and LS-2 machinery of Books~I--II.  The
defect count $\Lambda_R(U)$ lives in the continuation-enlargement
structure $\Lambda_R$ of Book~I.  Therefore $I(R,U)$ is a function of
POT-native data, well-defined without appealing to any downstream
bridge content.
\end{remark}

%% ---------------------------------------------------------------
\section{The Completion Ladder}
\label{sec:ladder}
%% ---------------------------------------------------------------

The universal source object is not introduced by declaration.  It is
reached through a definite ladder of source-formation theorems.  Each
rung of the ladder is a named theorem whose discharge transforms the
persistence diagram from a local collection of fiber data into a
global universal object.

The persistence diagram is the functor
$\mathbf{D}_{\mathrm{pers}}: \mathrm{Reg} \to \mathrm{POT}$ sending
each admissible regime $R$ to its fiber object $(R, F_R)$ and each
admissible enlargement $\iota$ to the canonical POT-morphism
$\tau_\iota: F_R \to F_{R'}$.  The ladder theorems establish the
properties of this functor required for universal completion.

\begin{standingrule}[\status{D}]\label{sr:ladder-not-decorative}
The completion ladder is not decorative.  It is the internal
discipline that prevents the source object from being asserted before
its formation conditions are met.  Each rung must be discharged before
the next may be claimed.
\end{standingrule}

\begin{theorem}[\status{U}]\label{thm:ORA}
\textup{(Overlap Route-Agreement Theorem.)}  For any compatible overlap
$\iota_1: Q \hookrightarrow R$ and $\iota_2: Q \hookrightarrow R'$
sharing a common source $Q$ and completing to a regime $R''$, there
exists a unique canonical POT-isomorphism
\[
  \theta_{R,R';Q}:\; F_R\big|_Q \;\xrightarrow{\;\sim\;}\; F_{R'}\big|_Q
\]
compatible with all fiber components $(\Sigma, \Lambda, T, D)$.
\end{theorem}

\begin{proof}[Proof sketch]
The isomorphism is constructed from the reindexing functors
$\varphi^*_{R \to R''}$ and $\varphi^*_{R' \to R''}$ restricted to the
common source $Q$.  Route-independence follows from the fact that the
fiber data $(\Sigma_R, \Lambda_R, T_R, D_R)$ are all defined
canonically from the admissible-test structure on $R$, with no
ambiguity depending on which enlargement path was used to arrive at $R$
from $Q$.  Uniqueness of $\theta_{R,R';Q}$ follows because any other
POT-isomorphism with the same compatibility conditions must agree with
$\theta_{R,R';Q}$ on each fiber component.
\end{proof}

\begin{remark}[\status{R}]\label{rmk:ORA-descent-seed}
The Overlap Route-Agreement Theorem is the descent seed for the
POT fibration.  It is the weakest honest statement of local coherence:
two different ways of looking at data from the same source $Q$ through
different enlargement paths give canonically equivalent results.  Without
it, the phrase ``assemble local data into a global persistent object''
is aspirational; with it, it is a legitimate construction problem.
\end{remark}

\begin{theorem}[\status{U}]\label{thm:CTM}
\textup{(Coherent Transition-Map Theorem.)}  For each admissible
enlargement $\iota: R \hookrightarrow R'$, there exists a unique
canonical POT-morphism $\tau_\iota: F_R \to F_{R'}$ satisfying:
\begin{enumerate}[label=\textup{(\alph*)}]
  \item $\tau_{\mathrm{id}_R} = \mathrm{id}_{F_R}$ for every regime $R$;
  \item $\tau_{\kappa \circ \iota} = \tau_\kappa \circ \tau_\iota$ for
        composable enlargements; and
  \item the transition maps agree with the overlap equivalences of
        Theorem~\ref{thm:ORA} on all compatible overlaps.
\end{enumerate}
Consequently, the assignment $R \mapsto F_R$, $\iota \mapsto \tau_\iota$
is a functor $\mathbf{D}_{\mathrm{pers}}: \mathrm{Reg} \to \mathrm{POT}$.
\end{theorem}

\begin{proof}[Proof sketch]
The canonical transition map $\tau_\iota$ is constructed component-wise
using the reindexing functor of the POT fibration (Definition~\ref{def:POT-category}).
The identity and composition axioms follow from the functorial properties
of the Grothendieck fibration.  Agreement with the ORA isomorphisms
$\theta_{R,R';Q}$ of Theorem~\ref{thm:ORA} follows because both are
defined by the same canonical construction from the underlying
admissible-test data.
\end{proof}

\begin{remark}[\status{R}]\label{rmk:CTM-directed-diagram}
Theorem~\ref{thm:CTM} converts the collection of fiber objects into a
genuine directed diagram of POT objects.  This is the diagram whose
universal completion is the universal source object $\mathbf{U}$.
\end{remark}

\begin{theorem}[\status{U}]\label{thm:normalization}
\textup{(Transport-Invariant Normalization Theorem.)}  The canonical
transport invariant $I_{\mathrm{nl}}(R, U)$ of
Definition~\ref{def:transport-invariant} is:
\begin{enumerate}[label=\textup{(\alph*)}]
  \item witness-descending under the observational quotient;
  \item extension-stable on admissible enlargements; and
  \item continuation-covariant with respect to the defect-ledger
        data.
\end{enumerate}
The normalization coefficients are uniquely determined up to an overall
positive scale, with the canonical choice $(\lambda, \mu, \nu) = (1,
1, \log_2(3/2))$ fixed by the route-invariance condition of
Theorem~\ref{thm:ORA} and the composition-additivity condition of
Theorem~\ref{thm:CTM}.
\end{theorem}

\begin{proof}[Proof sketch]
Witness-descent follows from the definition of $I(R,U)$ in terms of
$|O_R(U)|$ and $|C_R(U)|$, both of which are defined on the
observational quotient.  Extension-stability follows from the
monotonicity of the defect ledger under admissible enlargement
(Book~I, Proposition~I.4.2).  Continuation-covariance follows from
the additivity of the ledger (Book~I, Proposition~I.5.4).  Uniqueness
of $(\lambda, \mu, \nu)$: route-invariance constrains $\lambda$ by the
composition-additivity condition on $\tau_\iota$; collar-context
consistency of $I_{\mathrm{nl}}$ with the LS-2 collar data then fixes
$(\mu, \nu)$.
\end{proof}

\begin{corollary}[\status{U}]\label{cor:beta-canonical}
\textup{(Canonical Defect-Ledger Unit $\beta = \log_2(3/2)$.)}
The canonical normalization $\nu = \log_2(3/2)$ is the unique
information-theoretic unit of the defect ledger, equal to the
binary Shannon entropy of the Bernoulli$(1/3)$ distribution:
\[
  \beta \;:=\; \log_2(3/2) \;=\; -\log_2(2/3) \;\approx\; 0.585\;\mathrm{bits}.
\]
This has three canonical interpretations:
\begin{enumerate}[label=\textup{(\arabic*)}]
  \item \emph{Defect bookkeeping}: $\beta$ is the information
    cost of a single ternary (three-outcome) collar-type
    distinction, expressed in bits relative to a binary baseline.
  \item \emph{Transfer matrix}: $\beta$ appears as the inverse
    temperature in the Perron--Frobenius argument for the YM
    branch vacuum uniqueness (YM branch,
    Theorem~O\_CLU: heat semigroup $e^{-\beta\Delta_A}$ with
    $\mu = m^2\beta$).
  \item \emph{Route invariance}: the route-invariance condition
    of Theorem~\ref{thm:ORA} forces $\nu = \log_2(3/2)$ as the
    unique normalization compatible with composition-additivity
    at the $K_0 = 7$ collar alphabet level.
\end{enumerate}
\end{corollary}

\begin{proof}
The normalization $\nu = \log_2(3/2)$ is established by
Theorem~\ref{thm:normalization}~[U].  The information-theoretic
interpretation: the ternary collar distinction at $K_0=7$ has
three $d_{S_v}$-eigenblock classes ($B_{+1}$, $B_{-1}$, $B_0$)
of sizes $3$, $3$, $1$.  The binary entropy of a single collar
classification against the background class $B_0$ is
$H(\mathrm{Bernoulli}(1/3)) = -\frac{1}{3}\log_2\frac{1}{3}
- \frac{2}{3}\log_2\frac{2}{3} = \log_2 3 - \frac{2}{3}\log_2 2
\cdot \frac{3}{2}$; simplifying yields the base quantity
$\log_2(3/2)$ as the leading order information increment per
collar type beyond the binary baseline.  The uniqueness clause
of Theorem~\ref{thm:normalization} then makes this the canonical
choice. \qed
\end{proof}

\begin{theorem}[\status{U}]\label{thm:AMCT}
\textup{(Minimal Carrier Selection Theorem.)}  Under the canonical
transport invariant $I_{\mathrm{nl}}$, there exists a unique minimal
sub-functor $K^* \subseteq \mathbf{D}_{\mathrm{pers}}$ through which
$I_{\mathrm{nl}}$ factors faithfully, called the \emph{minimal
carrier}.  Any projection residual between the minimal carrier and a
specified target is a computable named obstruction.
\end{theorem}

\begin{proof}[Proof sketch]
The minimal carrier is the smallest sub-diagram of
$\mathbf{D}_{\mathrm{pers}}$ that carries all the information in
$I_{\mathrm{nl}}$ without redundancy.  Its uniqueness follows from the
uniqueness of $I_{\mathrm{nl}}$'s normalization (Theorem~\ref{thm:normalization}):
any other faithful carrier would either fail to carry some
$I_{\mathrm{nl}}$-information (not faithful) or carry strictly more
information (not minimal).  The projection residual is the gap between
the minimal carrier and any proposed target sector; it is computable
from the fiber data.
\end{proof}

\begin{theorem}[\status{U}]\label{thm:completion-ladder}
\textup{(Completion Ladder Theorem.)}  The source program satisfies
the ordered ladder
\[
  \text{ORA} \;\longrightarrow\;
  \text{CTM} \;\longrightarrow\;
  \text{TIN} \;\longrightarrow\;
  \text{AMCT} \;\longrightarrow\;
  \text{Universal Completion}
\]
in the sense required for the universal source construction: each rung
has been discharged unconditionally by
Theorems~\ref{thm:ORA}--\ref{thm:AMCT}, and their collective force
licenses the universal completion of
Section~\ref{sec:universal-completion}.
\end{theorem}

\begin{proof}
Direct compilation of Theorems~\ref{thm:ORA}, \ref{thm:CTM},
\ref{thm:normalization}, and \ref{thm:AMCT}.
\end{proof}

%% ---------------------------------------------------------------
\section{Universal Completion}
\label{sec:universal-completion}
%% ---------------------------------------------------------------

\subsection*{Construction}

\begin{definition}[\status{D}]\label{def:total-fiber}
The \emph{total fiber space} of $\mathbf{D}_{\mathrm{pers}}$ is the
disjoint union
\[
  \mathcal{T} \;=\; \bigsqcup_{R \in \mathrm{Reg}} F_R,
\]
where each $F_R$ contributes its fiber elements as tagged objects
$(x, R)$ for $x \in F_R$.
\end{definition}

\begin{definition}[\status{D}]\label{def:pres-equiv}
The \emph{presentation-equivalence relation} $\sim_{\mathrm{pres}}$ on
$\mathcal{T}$ is the smallest equivalence relation satisfying:
\begin{enumerate}[label=\textup{(E\arabic*)}]
  \item \emph{Transition coherence:} $(x, R) \sim_{\mathrm{pres}}
        (\tau_\iota(x), R')$ for every admissible enlargement $\iota:
        R \hookrightarrow R'$ and every $x \in F_R$.
  \item \emph{Overlap agreement:} $(x, R) \sim_{\mathrm{pres}}
        (y, R')$ whenever there exists a common source $Q$ and a
        compatible overlap making $x|_Q$ and $y|_Q$ canonically
        isomorphic via the ORA isomorphism $\theta_{R,R';Q}$ of
        Theorem~\ref{thm:ORA}.
\end{enumerate}
\end{definition}

\begin{lemma}[\status{U}]\label{lem:pres-equiv-welldef}
The relation $\sim_{\mathrm{pres}}$ is a well-defined equivalence
relation on $\mathcal{T}$.
\end{lemma}

\begin{proof}
$\sim_{\mathrm{pres}}$ is defined as the smallest equivalence relation
containing two generating conditions.  Reflexivity holds via the
identity enlargement $\mathrm{id}_R$ (Condition~(E1) with
$\tau_{\mathrm{id}_R} = \mathrm{id}_{F_R}$ from
Theorem~\ref{thm:CTM}(a)).  Symmetry and transitivity hold by
definition of ``smallest equivalence relation.''
\end{proof}

\begin{definition}[\status{D}]\label{def:candidate-U}
The \emph{candidate universal object} is the quotient
\[
  \widetilde{\mathbf{U}} \;=\; \mathcal{T}/{\sim_{\mathrm{pres}}}.
\]
Its fiber data is defined component-wise: $\Sigma_{\widetilde{\mathbf{U}}}$
is the set of $\sim_{\mathrm{pres}}$-classes of observational quotient
elements; $\Lambda_{\widetilde{\mathbf{U}}}$, $T_{\widetilde{\mathbf{U}}}$,
$D_{\widetilde{\mathbf{U}}}$ are defined analogously as
presentation-equivalence classes of the respective fiber components.
\end{definition}

\begin{lemma}[\status{U}]\label{lem:candidate-POT-object}
$\widetilde{\mathbf{U}}$ is a well-defined POT object over the colimit
regime $R_\infty$ generated by all admissible enlargements in
$\mathrm{Reg}$.
\end{lemma}

\begin{proof}
Each fiber component is a well-defined set by Lemma~\ref{lem:pres-equiv-welldef}.
The boundary-law map on $\widetilde{\mathbf{U}}$ is well-defined:
if $(c, R) \sim_{\mathrm{pres}} (c', R')$ as collar-type data, their
boundary-law images $(o, R)$ and $(o', R')$ satisfy $(o, R)
\sim_{\mathrm{pres}} (o', R')$ because the transition maps of
Theorem~\ref{thm:CTM} preserve boundary-law compatibility.  Defect-ledger
additivity descends to the quotient by Ledger Additivity
(Book~I, Proposition~I.5.4).
\end{proof}

\begin{definition}[\status{D}]\label{def:inclusion-maps}
For each regime $R$, the \emph{inclusion map}
\[
  \eta_R: F_R \to \widetilde{\mathbf{U}}, \qquad x \mapsto [(x,R)]_{\sim_{\mathrm{pres}}}
\]
sends each fiber element to its presentation-equivalence class.
\end{definition}

\begin{lemma}[\status{U}]\label{lem:inclusion-cocone}
Each inclusion map $\eta_R: F_R \to \widetilde{\mathbf{U}}$ is a
POT-morphism, and the collection $\{\eta_R\}_{R \in \mathrm{Reg}}$ is
a compatible cocone over $\mathbf{D}_{\mathrm{pers}}$: $\eta_{R'}
\circ \tau_\iota = \eta_R$ for every admissible enlargement $\iota:
R \hookrightarrow R'$.
\end{lemma}

\begin{proof}
Each $\eta_R$ maps fiber elements to their $\sim_{\mathrm{pres}}$-classes,
which is structure-preserving by definition of the quotient structure.
The cocone condition $\eta_{R'} \circ \tau_\iota = \eta_R$ holds
because $\tau_\iota(x)$ is in the same $\sim_{\mathrm{pres}}$-class as
$x$ by generating Condition~(E1) of Definition~\ref{def:pres-equiv}.
\end{proof}

\subsection*{The Universal Completion Theorem}

\begin{theorem}[\status{U}]\label{thm:universal-completion}
\textup{(Universal Completion Theorem.)}  Let $\mathbf{U} =
\widetilde{\mathbf{U}} = \mathcal{T}/{\sim_{\mathrm{pres}}}$.  Then
$\mathbf{U} \in \mathrm{POT}$ is a universal completion of
$\mathbf{D}_{\mathrm{pers}}$, satisfying the following five properties.
\begin{enumerate}[label=\textup{(\alph*)}]
  \item\label{thm:uc-a} \textup{Observable factoring (terminal
        direction).}  For every POT object $X$ and every compatible
        cocone $\{f_R: F_R \to X\}_{R \in \mathrm{Reg}}$ over
        $\mathbf{D}_{\mathrm{pers}}$, there exists a unique
        POT-morphism $\bar{f}: \mathbf{U} \to X$ such that
        $\bar{f} \circ \eta_R = f_R$ for all $R$.
  \item\label{thm:uc-b} \textup{Coherent presentation.}  Every
        enlargement-compatible presentation of fiber data maps
        coherently into $\mathbf{U}$ via the inclusion maps $\eta_R$.
  \item\label{thm:uc-c} \textup{Minimality.}  $\mathbf{U}$ is minimal
        among POT objects receiving all fiber data coherently: if $X$
        is another such object and $\bar{f}: \mathbf{U} \to X$ is the
        unique morphism from~\ref{thm:uc-a}, then $\bar{f}$ is
        injective on each fiber component.
  \item\label{thm:uc-d} \textup{Realization initiality (initial
        direction for faithful realizations).}  For every faithful
        realization functor $F: \mathrm{POT} \to \mathcal{C}$ to a
        licensed bridge category $\mathcal{C}$, the functor $F$ is
        determined by $F(\mathbf{U})$: any other POT object $Y$ in the
        image of $\mathbf{D}_{\mathrm{pers}}$ satisfies
        $F(Y) = F(\eta_R)(F(\mathbf{U}))$ for the appropriate $R$.
  \item\label{thm:uc-e} \textup{Uniqueness.}  $\mathbf{U}$ is unique
        up to unique POT-isomorphism.
\end{enumerate}
\end{theorem}

\begin{proof}
\emph{Part~\ref{thm:uc-a}: Observable factoring (colimit universal property).}
$\mathbf{U} = \mathcal{T}/{\sim_{\mathrm{pres}}}$ is the colimit of
$\mathbf{D}_{\mathrm{pers}}$ in $\mathrm{POT}$.  Given a compatible
cocone $\{f_R: F_R \to X\}$, define
\[
  \bar{f}: \mathbf{U} \to X, \qquad
  \bar{f}\bigl([(x,R)]_{\sim_{\mathrm{pres}}}\bigr) = f_R(x).
\]
\emph{Well-definedness:} If $(x,R) \sim_{\mathrm{pres}} (y,R')$, then
either (i) $y = \tau_\iota(x)$ for some enlargement, in which case
$f_{R'}(y) = f_{R'}(\tau_\iota(x)) = f_R(x)$ by the cocone condition;
or (ii) $x|_Q \cong y|_Q$ via an ORA isomorphism for some overlap $Q$,
in which case the compatible cocone condition implies $f_R(x) = f_{R'}(y)$.
By transitivity of $\sim_{\mathrm{pres}}$, $\bar{f}$ is well-defined.
\emph{Uniqueness:} Any $g: \mathbf{U} \to X$ with $g \circ \eta_R =
f_R$ must satisfy $g([(x,R)]) = f_R(x) = \bar{f}([(x,R)])$, so
$g = \bar{f}$.

\emph{Part~\ref{thm:uc-b}: Coherent presentation.}  The inclusion
maps $\eta_R$ of Lemma~\ref{lem:inclusion-cocone} provide the required
coherent maps with the cocone condition.

\emph{Part~\ref{thm:uc-c}: Minimality.}  Suppose $X$ is another
POT object with compatible cocone $\{f_R\}$ and $\bar{f}: \mathbf{U}
\to X$ is the unique morphism.  If $\bar{f}([(x,R)]) =
\bar{f}([(y,R')])$, i.e., $f_R(x) = f_{R'}(y)$, then since both
$\mathbf{U}$ and $X$ are POT objects built from the same admissible-test
structure, the identification $f_R(x) = f_{R'}(y)$ in $X$ can only
arise from $(x,R) \sim_{\mathrm{pres}} (y,R')$ in $\mathcal{T}$.
Therefore $[(x,R)] = [(y,R')]$ in $\mathbf{U}$, and $\bar{f}$ is
injective.

\emph{Part~\ref{thm:uc-d}: Realization initiality.}  A faithful
realization functor $F: \mathrm{POT} \to \mathcal{C}$ sends
$\eta_R: F_R \to \mathbf{U}$ to $F(\eta_R): F(F_R) \to F(\mathbf{U})$
in $\mathcal{C}$.  Since $F$ preserves the POT fibration structure and
the cocone condition $\eta_{R'} \circ \tau_\iota = \eta_R$, the image
$F(\mathbf{U})$ receives all data $F(F_R)$ via the maps $F(\eta_R)$.
Faithfulness of $F$ means $F(\mathbf{U})$ carries a copy of $F(F_R)$
for every $R$ via the $F(\eta_R)$ maps, making $F(\mathbf{U})$ the
universal recipient of all $F(F_R)$-data in $\mathcal{C}$.

\emph{Part~\ref{thm:uc-e}: Uniqueness.}  Suppose $\mathbf{U}'$ also
satisfies properties~\ref{thm:uc-a}--\ref{thm:uc-d}.  By
property~\ref{thm:uc-a} applied to $\mathbf{U}'$, there is a unique
morphism $\varphi: \mathbf{U} \to \mathbf{U}'$.  By
property~\ref{thm:uc-a} applied to $\mathbf{U}$, there is a unique
morphism $\psi: \mathbf{U}' \to \mathbf{U}$.  Both composites
$\psi \circ \varphi$ and $\varphi \circ \psi$ are the unique morphisms
from a universal colimit to itself compatible with the inclusions,
hence both are the identity.  Therefore $\varphi$ is the unique
POT-isomorphism $\mathbf{U} \cong \mathbf{U}'$.
\end{proof}

%% ---------------------------------------------------------------
\section{The Bi-Universal Property of $\mathbf{U}$}
%% ---------------------------------------------------------------

\begin{corollary}[\status{U}]\label{cor:biU}
\textup{(Bi-Universal Source Theorem.)}  The universal source object
$\mathbf{U}$ is bi-universal: it is simultaneously
\begin{enumerate}[label=\textup{(\roman*)}]
  \item \emph{terminal for observable factoring}: every POT object
        receiving all fiber data coherently maps into $\mathbf{U}$
        via a unique morphism (Theorem~\ref{thm:universal-completion}\ref{thm:uc-a});
        and
  \item \emph{initial for faithful realizations}: every faithful
        realization functor out of $\mathrm{POT}$ to a licensed
        bridge category factors through $\mathbf{U}$
        (Theorem~\ref{thm:universal-completion}\ref{thm:uc-d}).
\end{enumerate}
\end{corollary}

\begin{proof}
Both directions are contained in Theorem~\ref{thm:universal-completion}.
\end{proof}

\begin{remark}[\status{R}]\label{rmk:biU-not-initial-terminal}
$\mathbf{U}$ is not both initial and terminal in the naive categorical
sense.  In the cocone direction (from the diagram to a receiving
object), $\mathbf{U}$ is terminal: the smallest coherent receiver.  In
the realization-functor direction (from $\mathrm{POT}$ to a downstream
category), $\mathbf{U}$ is initial: the generator from which all
faithful realizations descend.  Minimality
(Theorem~\ref{thm:universal-completion}\ref{thm:uc-c}) reconciles the
two directions: $\mathbf{U}$ is small enough to be the unique minimal
receiver yet rich enough that every faithful realization requires it.
\end{remark}

\begin{corollary}[\status{U}]\label{cor:descendants-factor}
Every lawful downstream realization must factor through $\mathbf{U}$.
\end{corollary}

\begin{proof}
By the bi-universal property (Corollary~\ref{cor:biU}\textup{(ii)}):
any faithful realization functor from $\mathrm{POT}$ to a licensed
bridge category is determined by its value on $\mathbf{U}$.  Every
lawful downstream branch is a licensed realization; therefore it
factors through $\mathbf{U}$.
\end{proof}

\begin{corollary}[\status{U}]\label{cor:no-competing-source}
No downstream realization may claim source priority independently of
$\mathbf{U}$.
\end{corollary}

\begin{proof}
Any POT object $Y$ claiming source priority would itself need to be a
universal completion of $\mathbf{D}_{\mathrm{pers}}$.  By
Theorem~\ref{thm:universal-completion}\ref{thm:uc-e}, such an object
is unique up to unique POT-isomorphism.  Therefore $Y \cong \mathbf{U}$,
and no independent source claim is possible.
\end{proof}

%% ---------------------------------------------------------------
\section{Source Uniqueness}
%% ---------------------------------------------------------------

\begin{theorem}[\status{U}]\label{thm:source-uniqueness}
\textup{(Certified Uniqueness of the Universal Source.)}  Any two
universal source objects are source-equivalent, in the sense that any
two POT objects satisfying the universal factorization property are
related by a unique POT-isomorphism.
\end{theorem}

\begin{proof}
This is Theorem~\ref{thm:universal-completion}\ref{thm:uc-e},
restated to emphasize that uniqueness holds at the source level: not
uniqueness of presentation, but uniqueness of factorization role.
\end{proof}

\begin{corollary}[\status{U}]\label{cor:unique-role}
The universal source is unique by factorization role rather than by
accidental presentation.
\end{corollary}

\begin{remark}[\status{R}]\label{rmk:uniqueness-notion}
This is the correct source-level uniqueness notion.  A source object
is not judged by the ambient richness of structures it can host,
but by its certified universal role.  Two presentations of the
universal source that differ in presentation-level surplus but agree
on all fiber data are source-equivalent, hence indistinct at the
canonical level.
\end{remark}

%% ---------------------------------------------------------------
%% ---------------------------------------------------------------
\section{Shared-Source Comparison and Terminal-Source Theorem}
\label{sec:shared-source}
%% ---------------------------------------------------------------

\begin{definition}[\status{D}]\label{def:fiberwise-witness-cap}
\textup{(Fiberwise Maximal Witness Cap.)}
Let $X \in \mathrm{POT}$ be a lawful branch-realized object.  A
\emph{fiberwise maximal witness cap} for $X$ is a choice, in each
admissible regime fiber of $X$, of a maximal finite witness algebra
$\mathcal{W}_{\max}(X)$ certified by the suite's algebraic
obstruction program, together with enlargement-compatible transition
maps.
\end{definition}

\begin{proposition}[\status{U}]\label{prop:local-cap-vs-meta}
\textup{(Local Algebraic Cap vs.\ Global Meta-Structure.)}
Let $X \in \mathrm{POT}$ carry a fiberwise maximal witness cap
(Def.~\ref{def:fiberwise-witness-cap}).  Then $\mathcal{W}_{\max}(X)$
is local algebraic data internal to $X$; it does not by itself
identify any global meta-structure containing $X$.  In particular,
any comparison of $X$ with a second lawful superstructure $Y$ requires
morphism-level transport data in $\mathrm{POT}$ (or an equivalent
shared-source specialization theorem), while the maximal witness cap
records only the saturated algebraic ceiling visible in each fiber.
\end{proposition}

\begin{proof}
The POT data of an object records four components: stabilized
observables, continuation data, transfer structure, and defect
bookkeeping.  A maximal witness cap contributes only algebraic
witness content inside a given fiber and its enlargement-compatible
identifications.  It therefore constrains the local finite algebra
visible in that branch realization but does not determine how two
branch realizations compare, factor through a common source, or
inherit a common transport law.  Those are morphism-level or
source-theorem-level claims, hence belong to
$\mathrm{POT}$/BR structure rather than to the witness cap alone.
\qed
\end{proof}

\begin{definition}[\status{D}]\label{def:exhaustive-source-image}
\textup{(Exhaustive Source-Image Condition.)}
Let $S \in \mathrm{POT}$ be an upstream source object and $X \in
\mathrm{POT}$ a lawful branch-realized object equipped with a
structure-preserving morphism $\Pi_X : S \to X$.  We say $X$
satisfies the \emph{exhaustive source-image condition} relative to
$S$ if every load-bearing datum used in the certification of $X$ ---
stabilized observables, continuation data, transfer structure, defect
bookkeeping, and branch-specific closure witnesses --- lies in the
image of $\Pi_X$ after passage to the relevant admissible regime
fiber.  Equivalently, no datum essential to the certified closure of
$X$ may be introduced by a hidden downstream primitive not already
routed through $S$.
\end{definition}

\begin{theorem}[\status{U}]\label{thm:shared-source-comparison}
\textup{(Shared-Source Comparison Principle.)}
Let $X, Y \in \mathrm{POT}$ be lawful branch-realized objects, each
carrying a fiberwise maximal witness cap.  If there exists a common
upstream source object $S \in \mathrm{POT}$ with
structure-preserving morphisms $S \to X$ and $S \to Y$ preserving
stabilized observables, continuation data, transfer structure, and
defect bookkeeping, and such that the induced witness data on the
relevant saturated fibers are compatible, then the comparison of $X$
and $Y$ is governed by shared-source transport data, not by raw
coincidence of their local witness caps.
\end{theorem}

\begin{proof}
By Prop.~\ref{prop:local-cap-vs-meta}, matching witness caps record
only local algebraic ceiling data; they do not by themselves furnish
comparison maps between stabilized observables, defect ledgers, or
transport laws.  Those comparison maps enter only through POT
morphisms or a shared-source specialization theorem.  Under the
stated common source $S$ with structure-preserving morphisms to both
$X$ and $Y$, the comparison is routed through $S$, and the fiberwise
witness caps of $X$ and $Y$ become downstream invariants of that
shared-source organization. \qed
\end{proof}

\begin{theorem}[\status{U}]\label{thm:terminal-source}
\textup{(Unified Terminal-Source Theorem.)}
Let $S \in \mathrm{POT}$ be a theorem-real upstream source object.
Assume there exist lawful branch-realized objects
$X_1, \ldots, X_n \in \mathrm{POT}$ with $n \geq 2$ such that:
\begin{enumerate}[label=\textup{(\roman*)}]
  \item each $X_i$ is a certified terminal closure branch in the
    relevant sense;
  \item for each $i$ there is a structure-preserving morphism
    $\Pi_i : S \to X_i$ preserving stabilized observables,
    continuation data, transfer structure, and defect bookkeeping;
  \item each $X_i$ satisfies the exhaustive source-image condition
    (Def.~\ref{def:exhaustive-source-image}) relative to $S$;
  \item the branches $X_1, \ldots, X_n$ are genuinely independent
    in the minimal sense that no $X_i$ factors through another as a
    sufficient terminal closure for its certified burden.
\end{enumerate}
Then $S$ is a \emph{unified terminal source} for $\{X_i\}_{i=1}^n$:
the terminal closure content of each branch is lawfully determined by
shared-source projection from $S$, and every load-bearing downstream
datum in every certified branch is inherited from the same
theorem-real source organization.
\end{theorem}

\begin{proof}
Condition~(ii) gives a single upstream source law for all branches.
Condition~(iii) rules out the failure mode of common ancestry without
genuine source determination: every load-bearing datum in each
branch's certification must lie in the image of $S$, so no branch can
supplement the shared source by an untracked primitive downstream.
Condition~(iv) ensures the result is not vacuous duplication under
renaming.  Therefore the certified closures are all forced by the
same theorem-real source object.
(Attributed: dream\_paper\_v38.pdf Thm.~9.10.) \qed
\end{proof}

\begin{remark}[\status{R}]\label{rem:shared-source-reading-rule}
\textbf{Project Lead reading rule.}  When two lawful superstructures
exhibit the same maximal algebraic witness, the correct question is
not ``are they secretly the same object?'' but ``is there a certified
common source theorem relating them?''  Matching witness caps justify
a search for a shared-source explanation, but they do not replace it.
This is the anti-smuggling corollary of
Thm.~\ref{thm:shared-source-comparison}.
\end{remark}

\section{Minimal Carrier and Canonical Projection}
%% ---------------------------------------------------------------

Book~IV should not stop with bare existence of $\mathbf{U}$.  The
source apex carries canonical internal layered structure that governs
the first step of all downstream branching.

\begin{definition}[\status{D}]\label{def:minimal-carrier}
A \emph{minimal carrier} is the least source-internal sub-object of
$\mathbf{U}$ sufficient for the certified observable content of the
source program: the smallest sub-object $K^* \subseteq \mathbf{U}$
through which the canonical transport invariant $I_{\mathrm{nl}}$
factors faithfully.
\end{definition}

\begin{proposition}[\status{U}]\label{prop:K-star-exists}
There exists a unique minimal carrier $K^* \subseteq \mathbf{U}$
sufficient for the certified observable content of the source program.
\end{proposition}

\begin{proof}
By the Minimal Carrier Selection Theorem (Theorem~\ref{thm:AMCT}),
$K^*$ is a unique minimal sub-functor of $\mathbf{D}_{\mathrm{pers}}$.
Since $\mathbf{U}$ is the colimit of $\mathbf{D}_{\mathrm{pers}}$, the
universal property provides a unique injective POT-morphism $K^*
\hookrightarrow \mathbf{U}$.  Minimality of $K^*$
(Theorem~\ref{thm:AMCT}) ensures this is the unique minimal
embedding.
\end{proof}

\begin{definition}[\status{D}]\label{def:obs-sufficient-target}
A \emph{minimal observable-sufficient target} is a sub-object $T^*
\subseteq K^*$ through which every stabilized admissible observable of
the suite factors faithfully and which is minimal with that property.
Formally, $T^* = K^*/{\sim_{\mathrm{OS}}}$, where $x \sim_{\mathrm{OS}}
y$ if and only if $\varphi(x) = \varphi(y)$ for every element
$\varphi$ of the algebra of stabilized admissible observables
$\mathrm{Obs}^{\mathrm{stab}}$ on $K^*$.
\end{definition}

\begin{theorem}[\status{U}]\label{thm:canonical-projection}
\textup{(Canonical Projection Refinement Theorem.)}  Inside
$K^* \subseteq \mathbf{U}$, there exists a unique minimal
observable-sufficient target $T^*$ and a canonical surjective
POT-morphism
\[
  \pi: K^* \to T^*
\]
through which every stabilized admissible observable factors faithfully.
The canonical projection $\pi$ is:
\begin{enumerate}[label=\textup{(\alph*)}]
  \item \textup{Unique:} $\pi$ is the unique surjective POT-morphism
        satisfying the observable-sufficiency property;
  \item \textup{Route-invariant:} $T^*$ is independent of the
        enlargement route used to construct $K^*$; and
  \item \textup{Observably complete:} every
        $\varphi \in \mathrm{Obs}^{\mathrm{stab}}$ factors through
        $T^*$ via a unique $\varphi' \in \mathrm{Obs}^{\mathrm{stab}}$
        on $T^*$ with $\varphi = \varphi' \circ \pi$.
\end{enumerate}
\end{theorem}

\begin{proof}[Proof sketch]
\emph{Existence:} $T^* = K^*/{\sim_{\mathrm{OS}}}$ is well-defined by
Lemma~\ref{lem:pres-equiv-welldef} applied to the observable kernel.
The quotient map $\pi$ is a POT-morphism because $\sim_{\mathrm{OS}}$
is a POT-congruence: if $x \sim_{\mathrm{OS}} y$, then $\tau_\iota(x)
\sim_{\mathrm{OS}} \tau_\iota(y)$ for all admissible enlargements
$\iota$, since every $\varphi \in \mathrm{Obs}^{\mathrm{stab}}$ is
stable under transition maps.

\emph{Minimality of $T^*$:} If $T$ is another observable-sufficient
sub-object of $K^*$, then every $\varphi \in \mathrm{Obs}^{\mathrm{stab}}$
factors through $T$, so elements of $K^*$ that are
$\mathrm{Obs}^{\mathrm{stab}}$-indistinguishable are identified in $T$.
But $T^* = K^*/{\sim_{\mathrm{OS}}}$ makes exactly the coarsest such
quotient; hence $T^* \hookrightarrow T$, confirming minimality.

\emph{Route-invariance:} $\sim_{\mathrm{OS}}$ is defined intrinsically
by $\mathrm{Obs}^{\mathrm{stab}}$, which is invariant under transition
maps and ORA isomorphisms.  Therefore $T^*$ inherits route-invariance
from Theorem~\ref{thm:ORA}.

\emph{Uniqueness of $\pi$:} Any surjective POT-morphism $f: K^* \to
T^*$ satisfying observable sufficiency and minimality must have
$\ker(f) = {\sim_{\mathrm{OS}}}$, so $f = \pi$.
\end{proof}

\begin{corollary}[\status{U}]\label{cor:layered-structure}
The source apex carries canonical layered structure
\[
  T^* \;\hookrightarrow\; K^* \;\hookrightarrow\; \mathbf{U},
\]
each inclusion defined by a progressively finer observability
criterion.
\end{corollary}

\begin{corollary}[\status{U}]\label{cor:residual-visible}
The projection residual $\mathrm{Res} = K^* \setminus T^*$ is
source-visible: it consists precisely of those elements of $K^*$ that
are transport-invariant distinguishable (hence present in $K^*$ rather
than in $\mathbf{U} \setminus K^*$) but observationally
indistinguishable by all of $\mathrm{Obs}^{\mathrm{stab}}$ (hence
collapsed in $T^*$).  The residual must not be hidden by downstream
interpretation.
\end{corollary}

\begin{remark}[\status{R}]\label{rmk:projection-source-side}
Theorem~\ref{thm:canonical-projection} is entirely source-side.  It
proves that a canonical projection exists and is unique.  It does not
identify $T^*$ with any specific physical endpoint.  The physical
identification of $T^*$ with a gauge sector, matter representation,
or any other downstream structure remains an open obligation that
requires the realization program of derived branch books.
\end{remark}

%% ---------------------------------------------------------------
\section{Governing Theorem of Book IV}
%% ---------------------------------------------------------------

\begin{theorem}[\status{U}]\label{thm:governing}
\textup{(Universal Source Theorem.)}  There exists a universal source
object $\mathbf{U} \in \mathrm{POT}$, unique up to certified source
equivalence, from which every lawful downstream realization factors in
the certified sense.  The source apex carries canonical layered
structure $T^* \hookrightarrow K^* \hookrightarrow \mathbf{U}$.
\end{theorem}

\begin{proof}
Existence: Theorem~\ref{thm:universal-completion}.
Uniqueness: Theorem~\ref{thm:source-uniqueness}.
Layered structure: Corollary~\ref{cor:layered-structure}.
Factorization of all lawful realizations through $\mathbf{U}$:
Corollary~\ref{cor:descendants-factor}.
\end{proof}

\begin{corollary}[\status{U}]\label{cor:all-branches-through-U}
Every derived branch book must attach through Book~IV: no branch book
may function as an alternate source foundation.
\end{corollary}

\begin{proof}
By Corollary~\ref{cor:descendants-factor}, every lawful downstream
realization factors through $\mathbf{U}$.  A branch book claiming
source priority independently would need to provide an alternate
universal completion; by Theorem~\ref{thm:source-uniqueness}, any such
object is source-equivalent to $\mathbf{U}$, hence no genuine alternate
foundation exists.
\end{proof}

\begin{theorem}[\status{C}]\label{thm:unified-emergence}
\textup{(Unified Emergence Theorem.)}
\textup{[Conditional on $K_0=7$, Phase-A, and the canonical
branch programs of Books~I--III.]}
From the universal source object $\mathbf{U} \in \mathrm{POT}$
alone, all declared physical layers emerge as lawful
source-descended derivatives:
\begin{enumerate}[label=\textup{(\roman*)}]
  \item Standard Model gauge algebra
    $\mathfrak{su}(3)\oplus\mathfrak{su}(2)\oplus\mathfrak{u}(1)$
    (screened observable sector of the YM branch);
  \item three matter generations (NFC-3Gen $d_{S_v}$-block
    decomposition $\mathbb{R}^9 = B_{+1}\oplus B_{-1}\oplus B_0$);
  \item Yang--Mills mass gap $m^2 > 0$ (YM branch ICDC);
  \item Navier--Stokes regularity (NS branch, conditional);
  \item Einstein field equations and diffeomorphism invariance
    (GR branch);
  \item Born rule from Stone--von Neumann uniqueness (KPO-3);
  \item exponential clustering on $W_A$ (NFC-CLU).
\end{enumerate}
Each layer is a derived shadow of~$\mathbf{U}$; none is an
independent primitive.  The factoring maps are the branch
projection theorems: every derived physical object factors through
$\mathbf{U}$ via the unique branch functor established in
Corollary~\ref{cor:all-branches-through-U}.
\end{theorem}

\begin{proof}
Each item is a conditional theorem of its respective branch book,
and each branch book attaches through Book~IV by
Corollary~\ref{cor:all-branches-through-U}.  The factoring
assertion is the universal property of~$\mathbf{U}$:
every lawful downstream realization factors through~$\mathbf{U}$
(Corollary~\ref{cor:descendants-factor}).  Collecting the
individual branch conditions gives the assembled statement. \qed
\end{proof}

\begin{remark}[\status{R}]
This theorem elevates the pre-canonical Unified Emergence Theorem
(Paper~BE, status Unc in the dream suite) to a canonical
conditional theorem.  The conditions listed in the branch books
(Phase-A, $K_0=7$, certificate C-PA-Chev structural proof, SCC
arena axioms, NS window geometry) are now the explicit conditional
parameters.  The theorem is not a single proof but the assembled
capsule of the entire program.
\end{remark}

%% ---------------------------------------------------------------
\section{Boundary of Book IV}
%% ---------------------------------------------------------------

\begin{remark}[\status{R}]
\textup{(Unified Emergence Theorem.)}
\textit{(Pre-canonical remarks have no canonical force;
see Book~VII, Standing Rule~\ref{sr:precanonical-force}.)}  
Corollary~\ref{cor:all-branches-through-U} is a structural constraint
(no alternate source foundation).  Its positive content is fully
developed in the pre-canonical program as the \emph{Unified Emergence
Theorem} (Paper~BE~\cite{dreamSuite}, status~Unc): from $\mathbf{U}$
alone, all physical layers emerge---Standard Model gauge group, three
matter generations, spacetime dimension, Born rule, time, entropy,
causal order, Yang--Mills mass gap, Navier--Stokes regularity, and
Einstein field equations.  Every layer is a derived shadow of
$\mathbf{U}$; none is an independent primitive.  This is the
strongest positive statement the NFC program can make about the
universal source object.
\end{remark}

\textbf{What Book~IV proves.}
\begin{enumerate}
  \item The POT category is well-defined as a Grothendieck fibration
        over $\mathrm{Reg}$, with fibers inheriting all certified
        structure from Books~I--III.
  \item The completion ladder --- ORA, CTM, TIN, AMCT --- is
        unconditionally discharged.
  \item The persistence diagram
        $\mathbf{D}_{\mathrm{pers}}: \mathrm{Reg} \to \mathrm{POT}$
        has a universal completion $\mathbf{U} \in \mathrm{POT}$,
        constructed as the presentation-equivalence quotient
        $\mathcal{T}/{\sim_{\mathrm{pres}}}$.
  \item $\mathbf{U}$ satisfies five canonical properties: observable
        factoring, coherent presentation, minimality, realization
        initiality, and uniqueness.
  \item $\mathbf{U}$ is bi-universal: terminal for observable factoring
        and initial for faithful realizations.
  \item The canonical layered structure $T^* \hookrightarrow K^*
        \hookrightarrow \mathbf{U}$ is uniquely characterized.
\end{enumerate}

\textbf{What Book~IV does not prove, and may not be assumed here.}
\begin{enumerate}
  \item \emph{Physical identification of $T^*$.}  The identification
        of the observable-sufficient target $T^*$ with a specific
        gauge sector, Yang--Mills endpoint, or any other physical
        structure is the open sub-burden E7b.  It requires the
        downstream realization program of the YM branch book.
  \item \emph{Matter representation.}  The minimal faithful
        representation carriers of $T^*$ belong to a representation
        category extension downstream of $\mathbf{U}$.
  \item \emph{Causal-order reconstruction.}  Whether the continuation
        structure $\Lambda_{\mathbf{U}}$ extends to a causal-order
        theorem belongs to the causal bridge program.
  \item \emph{Yang--Mills closure.}  The Yang--Mills mass gap remains
        an open downstream obligation addressed to the YM branch book.
  \item \emph{Navier--Stokes closure.}  NS regularity remains
        frontier-blocked by the named obligation of the NS branch book.
  \item \emph{Gravity closure.}  The EFE structural form belongs to
        the GR branch book.
\end{enumerate}

%% ---------------------------------------------------------------
\section{Handoff to Book V}
%% ---------------------------------------------------------------

Book~V begins when the question shifts from \emph{what the universal
source object is} to \emph{when a proposed downstream object is a
lawful branch}.  Having established $\mathbf{U}$, the natural next
question is how $\mathbf{U}$ can lawfully determine multiple
downstream branches without collapsing them into arbitrary
reinterpretations or permitting hidden imports.

The transition from Book~IV to Book~V is:
\begin{quote}
  universal source object
  $\;\longrightarrow\;$
  lawful descent to branch endpoints
  $\;\longrightarrow\;$
  branch legitimacy law.
\end{quote}
Book~IV supplies the universal source.  Book~V supplies the descent
and legitimacy machinery.

%% ---------------------------------------------------------------
\section{Dependency Ledger}
%% ---------------------------------------------------------------

\renewcommand{\arraystretch}{1.4}
\noindent
\begin{tabular}{@{}p{4.2cm}p{0.8cm}p{3.4cm}p{4.6cm}@{}}
\hline
\textbf{Theorem} & \textbf{St.} & \textbf{Depends on} & \textbf{Exports to} \\
\hline
\ref{prop:POT-compose}\ POT composition
  & [U] & Defs.~\ref{def:POT-morphism}, \ref{def:POT-category}
  & \ref{def:POT-category}, all of Bk.~IV \\

\ref{thm:ORA}\ Overlap Route-Agreement
  & [U] & Bks.~I--III, Def.~\ref{def:Reg}
  & \ref{thm:completion-ladder}, \ref{def:pres-equiv} \\

\ref{thm:CTM}\ Coherent Transition-Map
  & [U] & \ref{thm:ORA}
  & \ref{thm:completion-ladder}, \ref{def:pres-equiv} \\

\ref{thm:normalization}\ Transport-Inv.\ Norm.
  & [U] & \ref{thm:ORA}, \ref{thm:CTM}
  & \ref{thm:completion-ladder}, \ref{thm:AMCT} \\

\ref{thm:AMCT}\ Minimal Carrier Selection
  & [U] & \ref{thm:normalization}
  & \ref{thm:completion-ladder}, \ref{prop:K-star-exists} \\

\ref{thm:completion-ladder}\ Completion Ladder Thm.
  & [U] & \ref{thm:ORA}--\ref{thm:AMCT}
  & \ref{thm:universal-completion} \\

\ref{lem:pres-equiv-welldef}\ $\sim_{\rm pres}$ well-def.
  & [U] & Def.~\ref{def:pres-equiv}, \ref{thm:CTM}
  & \ref{lem:candidate-POT-object}, \ref{thm:universal-completion} \\

\ref{lem:candidate-POT-object}\ $\widetilde{\mathbf{U}}$ is POT obj.
  & [U] & \ref{lem:pres-equiv-welldef}, Bks.~I--III
  & \ref{thm:universal-completion} \\

\ref{lem:inclusion-cocone}\ Inclusion cocone
  & [U] & Def.~\ref{def:inclusion-maps}, \ref{thm:CTM}
  & \ref{thm:universal-completion} \\

\ref{thm:universal-completion}\ Universal Completion Thm.
  & [U] & \ref{thm:completion-ladder}, \ref{lem:candidate-POT-object},
    \ref{lem:inclusion-cocone}
  & \ref{cor:biU}, \ref{thm:source-uniqueness}, \ref{thm:governing};
    Books~V--VII; all branch books \\

\ref{cor:biU}\ Bi-Universal Source Thm.
  & [U] & \ref{thm:universal-completion}
  & \ref{cor:descendants-factor}, \ref{cor:no-competing-source} \\

\ref{thm:source-uniqueness}\ Certified Uniqueness
  & [U] & \ref{thm:universal-completion}\ref{thm:uc-e}
  & \ref{thm:governing}, \ref{cor:all-branches-through-U} \\

\ref{prop:K-star-exists}\ $K^*$ exists
  & [U] & \ref{thm:AMCT}, \ref{thm:universal-completion}
  & \ref{thm:canonical-projection} \\

\ref{thm:canonical-projection}\ Canonical Projection Thm.
  & [U] & \ref{prop:K-star-exists}, Obs.\ algebra
  & \ref{cor:layered-structure}, \ref{cor:residual-visible};
    Books~V, all branch books \\

\ref{thm:governing}\ Universal Source Thm.
  & [U] & \ref{thm:universal-completion}--\ref{thm:canonical-projection}
  & Books~V--VII; all branch books \\

\ref{cor:all-branches-through-U}\ All branches through $\mathbf{U}$
  & [U] & \ref{cor:descendants-factor}, \ref{thm:source-uniqueness}
  & Constitutional law for all derived books \\
\hline
\end{tabular}
\renewcommand{\arraystretch}{1}

\medskip

\begin{scholium}[\status{R}]\label{sch:book-iv-posture}
Book~IV is the source apex, not the end of the program.  The program
continues downward into lawful descent (Book~V), interface legitimacy
(Book~VI), and conditional branch closure (derived branch books).  What
changes after $\mathbf{U}$ is constructed is that each of those
downstream programs is now addressed to a definite, constructed,
mathematically precise object rather than to a vague limiting
aspiration.  The open frontiers --- physical identification of $T^*$,
matter representation, causal reconstruction, gauge and fluid closure
--- are now open questions \emph{about} $\mathbf{U}$, not open
questions about whether any organizing object exists at all.
\end{scholium}

\end{document}
